% GENERATED FILE. Edit canonical lessons, then run npm run generate:books.
% canonical-source-hash: fnv1a64:063f8027d9efb62f
% canonical-lessons: LA-C521-aegrotus, LA-C521-aegroto, LA-C521-languidus, LA-C521-debilis, LA-C521-debilitas

\chapter{Being ill}
\label{ch:la-being-ill}
\hlchaptermodality{521}

\begin{chapteropening}
\textbf{By the end of this chapter:} \emph{I can say that someone is ill and weak.}

Complete the last lesson of chapter 521: I can say that someone is ill and weak.
\end{chapteropening}

\section[aegrōtus, aegrōta, aegrōtum]{aegrōtus, aegrōta, aegrōtum — sick, ill}

\label{lesson:LA-C521-aegrotus}



\begin{quote}
Before the new one: say the Latin for to yawn, then the Latin for to sneeze.
\end{quote}

\subsection*{You'll want to know: aegrōtus, aegrōta, aegrōtum}
\textbf{aegrōtus, aegrōta, aegrōtum} — \textquotedblleft{}sick, ill\textquotedblright{}.

\begin{quote}
Fīlius aegrōtus in lectō iacet.
\end{quote}

\begin{quote}
\emph{The sick son lies in bed.}
\end{quote}

Built on \emph{aeger}. As a noun, \emph{aegrōtus} — a patient.

\begin{cousinweb}
\textbf{aegrōtus} is built on \emph{aeger} \textquotedblleft{}sick\textquotedblright{} (Lewis \& Short; EtymDB).
\end{cousinweb}

\begin{grammarlens}[title={What you've built}]
One more word for telling, explaining and describing at B1.
\end{grammarlens}

\subsection*{Your turn}
\begin{itemize}
\raggedright
  \item \emph{Say it:} \emph{aegrōtus, aegrōta, aegrōtum}
  \item \emph{Say it:} \emph{aegrōtus, aegrōta, aegrōtum}, once more
  \item \emph{Say it:} a sentence of your own with \emph{aegrōtus, aegrōta, aegrōtum}
  \item \emph{Recall:} say the Latin for to yawn, then the Latin for to sneeze, then say \emph{aegrōtus, aegrōta, aegrōtum} again
\end{itemize}

\begin{tcolorbox}[breakable,colback=teal!4,colframe=teal!35!black,title={Before you move on}]
What does aegrōtus, aegrōta, aegrōtum mean? (\textquotedblleft{}Sick, ill\textquotedblright{}.) Say it once more.
\end{tcolorbox}
\section[aegrōtō, aegrōtāre]{aegrōtō, aegrōtāre — to be ill}

\label{lesson:LA-C521-aegroto}



\begin{quote}
Before the new one: say the Latin for to sneeze, then the Latin for sick.
\end{quote}

\subsection*{You'll want to know: aegrōtō, aegrōtāre}
\textbf{aegrōtō, aegrōtāre} — \textquotedblleft{}to be ill\textquotedblright{}.

\begin{quote}
Avus iam trēs diēs aegrōtat.
\end{quote}

\begin{quote}
\emph{Grandfather has been ill for three days now.}
\end{quote}

\emph{aegrōtō, aegrōtāre, aegrōtāvī}, from \emph{aegrōtus}. With \emph{iam}, the present means has been.

\begin{cousinweb}
\textbf{aegrōtō} is built on \emph{aegrōtus} \textquotedblleft{}sick\textquotedblright{} (Lewis \& Short).
\end{cousinweb}

\begin{grammarlens}[title={What you've built}]
One more word for telling, explaining and describing at B1.
\end{grammarlens}

\subsection*{Your turn}
\begin{itemize}
\raggedright
  \item \emph{Say it:} \emph{aegrōtō, aegrōtāre}
  \item \emph{Say it:} \emph{aegrōtō, aegrōtāre}, once more
  \item \emph{Say it:} a sentence of your own with \emph{aegrōtō, aegrōtāre}
  \item \emph{Recall:} say the Latin for to sneeze, then the Latin for sick, then say \emph{aegrōtō, aegrōtāre} again
\end{itemize}

\begin{tcolorbox}[breakable,colback=teal!4,colframe=teal!35!black,title={Before you move on}]
What does aegrōtō, aegrōtāre mean? (\textquotedblleft{}To be ill\textquotedblright{}.) Say it once more.
\end{tcolorbox}
\section[languidus, languida, languidum]{languidus, languida, languidum — faint, weak, listless}

\label{lesson:LA-C521-languidus}



\begin{quote}
Before the new one: say the Latin for sick, then the Latin for to be ill.
\end{quote}

\subsection*{You'll want to know: languidus, languida, languidum}
\textbf{languidus, languida, languidum} — \textquotedblleft{}faint, weak, listless\textquotedblright{}.

\begin{quote}
Post febrem corpus languidum est.
\end{quote}

\begin{quote}
\emph{After the fever the body is weak.}
\end{quote}

From \emph{langueō} — to be faint. English \emph{languid}.

\begin{cousinweb}
\textbf{languidus} is built on \emph{langueō} \textquotedblleft{}to be faint, weary\textquotedblright{} (Lewis \& Short; EtymDB). English \emph{languid} comes from it, and \emph{languish} from the same verb.
\end{cousinweb}

\begin{grammarlens}[title={What you've built}]
One more word for telling, explaining and describing at B1.
\end{grammarlens}

\subsection*{Your turn}
\begin{itemize}
\raggedright
  \item \emph{Say it:} \emph{languidus, languida, languidum}
  \item \emph{Say it:} \emph{languidus, languida, languidum}, once more
  \item \emph{Say it:} a sentence of your own with \emph{languidus, languida, languidum}
  \item \emph{Recall:} say the Latin for sick, then the Latin for to be ill, then say \emph{languidus, languida, languidum} again
\end{itemize}

\begin{tcolorbox}[breakable,colback=teal!4,colframe=teal!35!black,title={Before you move on}]
What does languidus, languida, languidum mean? (\textquotedblleft{}Faint, weak, listless\textquotedblright{}.) Say it once more.
\end{tcolorbox}
\section[dēbilis, dēbile]{dēbilis, dēbile — weak, disabled}

\label{lesson:LA-C521-debilis}



\begin{quote}
Before the new one: say the Latin for to be ill, then the Latin for faint.
\end{quote}

\subsection*{You'll want to know: dēbilis, dēbile}
\textbf{dēbilis, dēbile} — \textquotedblleft{}weak, disabled\textquotedblright{}.

\begin{quote}
Senex dēbilis lentē ambulat.
\end{quote}

\begin{quote}
\emph{The weak old man walks slowly.}
\end{quote}

Two endings. Of a limb that will not work, or of general weakness.

\begin{cousinweb}
\textbf{dēbilis} is built on \emph{dē-} \textquotedblleft{}away from\textquotedblright{} and \emph{habilis} \textquotedblleft{}handy, able\textquotedblright{} (Lewis \& Short; EtymDB).
\end{cousinweb}

\begin{grammarlens}[title={What you've built}]
One more word for telling, explaining and describing at B1.
\end{grammarlens}

\subsection*{Your turn}
\begin{itemize}
\raggedright
  \item \emph{Say it:} \emph{dēbilis, dēbile}
  \item \emph{Say it:} \emph{dēbilis, dēbile}, once more
  \item \emph{Say it:} a sentence of your own with \emph{dēbilis, dēbile}
  \item \emph{Recall:} say the Latin for to be ill, then the Latin for faint, then say \emph{dēbilis, dēbile} again
\end{itemize}

\begin{tcolorbox}[breakable,colback=teal!4,colframe=teal!35!black,title={Before you move on}]
What does dēbilis, dēbile mean? (\textquotedblleft{}Weak, disabled\textquotedblright{}.) Say it once more.
\end{tcolorbox}
\section[dēbilitās, dēbilitātis]{dēbilitās, dēbilitātis — weakness, infirmity}

\label{lesson:LA-C521-debilitas}



\begin{quote}
Before the new one: say the Latin for faint, then the Latin for weak.
\end{quote}

\subsection*{You'll want to know: dēbilitās, dēbilitātis}
\textbf{dēbilitās, dēbilitātis} — \textquotedblleft{}weakness, infirmity\textquotedblright{}.

\begin{quote}
Dēbilitās corporis eum domī tenet.
\end{quote}

\begin{quote}
\emph{Weakness of body keeps him at home.}
\end{quote}

Feminine, third declension, from \emph{dēbilis}. English \emph{debility}.

\begin{cousinweb}
\textbf{dēbilitās} is built on \emph{dēbilis} \textquotedblleft{}weak\textquotedblright{} (Lewis \& Short; EtymDB). English \emph{debility} comes from it.
\end{cousinweb}

\begin{grammarlens}[title={What you've built}]
One more word for telling, explaining and describing at B1.
\end{grammarlens}

\subsection*{Your turn}
\begin{itemize}
\raggedright
  \item \emph{Say it:} \emph{dēbilitās, dēbilitātis}
  \item \emph{Say it:} \emph{dēbilitās, dēbilitātis}, once more
  \item \emph{Say it:} a sentence of your own with \emph{dēbilitās, dēbilitātis}
  \item \emph{Recall:} say the Latin for faint, then the Latin for weak, then say \emph{dēbilitās, dēbilitātis} again
\end{itemize}

\begin{tcolorbox}[breakable,colback=teal!4,colframe=teal!35!black,title={Before you move on}]
What does dēbilitās, dēbilitātis mean? (\textquotedblleft{}Weakness, infirmity\textquotedblright{}.) Say it once more.
\end{tcolorbox}
